\documentclass{article}

\usepackage{amsmath,amssymb}
\usepackage{xurl}
\usepackage[hidelinks]{hyperref}
\setlength{\emergencystretch}{2em}

% Shared commands for notes in docs/paper-gaps/.

\DeclareUnicodeCharacter{2080}{\ifmmode{}_{0}\else\textsubscript{0}\fi}
\DeclareUnicodeCharacter{2081}{\ifmmode{}_{1}\else\textsubscript{1}\fi}
\DeclareUnicodeCharacter{2082}{\ifmmode{}_{2}\else\textsubscript{2}\fi}
\DeclareUnicodeCharacter{2099}{\ifmmode{}_{n}\else\textsubscript{n}\fi}

\newcommand{\Lean}{\textsc{Lean}}
\ExplSyntaxOn
\NewDocumentCommand{\leanid}{m}{
  \ifmmode
    \textsf{\detokenize{#1}}
  \else
    \group_begin:
    \urlstyle{sf}
    \tl_set:Nn \l_tmpa_tl {#1}
    \tl_replace_all:Nnn \l_tmpa_tl {\_} {_}
    \exp_args:NV \path \l_tmpa_tl
    \group_end:
  \fi
}
\ExplSyntaxOff
\providecommand{\tr}{\operatorname{tr}}
\newcommand{\tnleanIssueBase}{https://github.com/LionSR/TNLean/issues/}
\newcommand{\tnleanPrBase}{https://github.com/LionSR/TNLean/pull/}
\newcommand{\tnleanDocBase}{https://sirui-lu.com/TNLean/docs/}
\newcommand{\ghissue}[1]{\href{\tnleanIssueBase#1}{GitHub issue \##1}}
\newcommand{\ghpr}[1]{\href{\tnleanPrBase#1}{PR \##1}}
\newcommand{\doclink}[1]{\href{\tnleanDocBase#1.html}{\path{#1}}}

% Dirac notation and the identity operator, as in blueprint/src/macros/common.tex.
% \providecommand keeps note-local definitions authoritative.
\usepackage{dsfont}
\providecommand{\ket}[1]{|#1\rangle}
\providecommand{\bra}[1]{\langle #1|}
\providecommand{\braket}[2]{\langle #1 | #2 \rangle}
\providecommand{\ketbra}[2]{|#1\rangle\!\langle #2|}
\providecommand{\Id}{\mathds{1}}
\providecommand{\one}{\mathds{1}}
\providecommand{\C}{\mathbb{C}}
\providecommand{\MN}[1]{M_{#1}(\C)}
\providecommand{\MD}{M_D(\C)}

% Verdict marker: \gapnote{<kind>}{<status>}. Kind is the policy
% classification (clarification, local-correction, scope-restriction,
% unfaithful, false-source, open-gap); status is open, wip (elimination
% underway), resolved, or historical. Machine-read by the site index;
% prints a verdict line.
\newcommand{\gapnote}[2]{\par\noindent\textbf{Verdict:} #1 (#2).\par}


% Fallbacks keep this note compilable when a different command.tex is found first.
\providecommand{\leanid}[1]{\textsf{\detokenize{#1}}}
\providecommand{\gapnote}[2]{\par\noindent\textbf{Verdict:} #1 (#2).\par}

\title{Input and Output Coefficients in the Scanner Estimate}
\author{}
\date{2026-10-16}

\begin{document}
\maketitle
\gapnote{local-fix}{coefficient convention retained}

\section{The source proposition and the coefficient convention}

Proposition~9.2 of the September~24, 2026 two-dimensional area-law
manuscript~\cite{OpenAI2026AreaLaw} (\path{08-scanner.tex}, lines~349--397)
bounds the replica defect energy produced by a collar scan at scale $n$. Let
$X$ be the target region, $Q_j$ the collar shells, $L$ the collar length, $m$
the band spacing, $K$ the number of bands, $D$ the charge depth, $r_0$ the
Hamiltonian truncation radius, $W$ the total weight with per-band weight
$a=W/K$, and $\ell,\kappa,\mu,\nu$ the fixed scanner exponents, with
$\nu>0$ and $\epsilon=n^{-\nu}$. Assuming the
inductive entropy bounds $S_\Omega(X)\le C_en^{1+e}$ and
$S_\Omega(Q_j)\le C_e(nL^e+n)$ for the ground state $\Omega$, Proposition~9.2
asserts that for every sufficiently large replica count $k$, some charge round
and some interpolation parameter
$p_k\in[\epsilon/2,\epsilon]$ satisfy the defect-energy bound
\[
    E_{\mathrm{def},k}\le Cn^\ell D^2W^2(\log n)^C
    \bigl(\delta_n^{1/8}+\epsilon+n^{-100}\bigr)+Cn^{-1000}+o_k(1),
\]
where $\delta_n=C_e\epsilon^{-1}(n^{e-\mu}+a^{1/4}(\log n)^C)$, together with
the comparator bounds $B_{\mathrm{sh}}\le C_e(nL^e+n)$,
$B_{\mathrm{exc}}\le CnD$, and $\mathcal B\le CnD(\log n)^C$.

The paper's proof of Proposition~9.2 (lines~399--533) combines five components
for the collar scan:
(1)~truncating the Hamiltonian near the scan region and bounding the resulting
trace-distance and regional-entropy perturbations (Proposition~4.5 and
Lemma~2.1);
(2)~constructing the random scan histories and proving splitting-ball sampling,
incidence dilution, bad-history tail bounds, and band nesting (Lemma~9.1);
(3)~applying the replica transport inequalities along each fill and charge
round (Proposition~7.4);
(4)~bounding the marginal log-moment generating functions, typical-window
masses, and initial/terminal comparator norms (Lemma~3.1 and
Proposition~8.1); and
(5)~telescoping the logarithmic norm over all $2nm$ rounds to select a charge
round and parameter $p_k$ with $\mathcal Q(p_k)/(KnD)\le\delta_n+\rho_k$ and
applying H\"older's inequality to bound $E_{\mathrm{def},k}$.

\paragraph{The scanner conclusion for the original physical system.}
The scanner theorem\footnote{The scanner conclusion:
    \leanid{TNLean.PEPS.AreaLaw.Scan.CollarScan.scannerStatement}.}
constructs a truncated ground vector, a target typical set and a common
Schmidt--Bell label sequence from the original gapped Hamiltonian, the
scanner geometry and the entropy hypotheses. The truncation coefficient
is chosen before the local dimension and the entropy coefficient. The
other uniform constants and the size threshold precede the finite
physical system and the total weight; the state and labels precede the
replica limit.

For every sufficiently large replica count, one charge round and one
parameter in $[\epsilon/2,\epsilon]$ satisfy the selected-density and
defect-energy estimates and both rough norm comparisons. Under the
additional sharp smallness condition, both sharp comparisons hold at
that same point. The marginal moment estimate applies to every region of
the same truncated vector, and one common marginal parameter controls
all status cuts. The final-round rough upper estimate has its own
vanishing replica remainder.

The conclusion does not assume a record of transport or norm-comparison
estimates. Those estimates are derived for the same physical data. The
abstract scanner deduction remains a separate result, while the direct
physical theorem uses the actual charge selection, the supplied-state
norm comparisons and the common status-cut bounds. The mathematical
adjustment recorded in this note concerns the output coefficients.

\paragraph{The entropy input constant and the enlarged output constants.}
The symbol $C_e$ is reused in the density and shell bounds of
\path{08-scanner.tex}, lines~379--389, after fixed factors and additive
errors are absorbed at lines~462--490. The proposition permits its
constants to depend on $C_e$ at line~396. To distinguish the two roles,
retain $C_e^{\rm in}$ in the original entropy hypotheses and put
$C_* = \max\{C_e^{\rm in},1\}$. This enlargement preserves those
hypotheses and introduces no assumption that $C_e^{\rm in}$ is positive.
The transferred whole-prefix entropy bound is
\begin{align}
    B_{\rm sh}
     & \le C_*(nL^e+n)+n^{-247}
    \le (C_*+1)(nL^e+n).\notag
\end{align}
The last inequality holds for $n\ge2$ and retains the entropy-continuity
error by enlarging the coefficient. The mismatch bound is
$B_{\rm exc}=1+3nD\log q\le(1+3\log q)nD$ when $D\ge1$.
A common output coefficient may therefore be chosen to dominate
$C_*+1$, $1+3\log q$ and both marginal-budget coefficients.

The density coefficient is likewise chosen after $C_e^{\rm in}$, but
before the finite physical system and total weight. It is a positive
constant $C_\delta$, giving
\begin{align}
    \delta_n
     & =C_\delta\epsilon^{-1}
    \left(n^{e-\mu}+a^{1/4}(\log n)^{C_{l,\delta}}\right).\notag
\end{align}
For example, the rough budget has coefficient
$C_{\rm rough}=6C_*+13+6\log q$, and the density conversion permits
$C_{\delta,0}=4(C_1+1)(C_{\rm rough}+2C_t)C_s$, where $C_t\ge1$
bounds the transport coefficient and $C_s\ge1$ bounds the rounded
scales. A subsequent enlargement serves both the density and energy
estimates. These constants absorb the interval length, typical-width,
mismatch and transport terms.

For $C_e^{\rm in}=0$, the output coefficient $C_\delta$ may remain
positive and the transferred shell entropy is not asserted to vanish.
The statement with each output coefficient numerically identical to
$C_e^{\rm in}$ is stronger than the coefficient convention used here;
this proof does not establish it. The original entropy hypotheses remain
unchanged.

\section{Step 1: Truncation trace distance and entropy continuity}

The paper states that the trace distance $t=\frac12\||\widetilde\Omega\rangle
    \langle\widetilde\Omega|-|\Omega\rangle\langle\Omega|\|_1$ between the original
and truncated ground states satisfies $t=O(n^{-500})$~\cite{OpenAI2026AreaLaw}
(\path{08-scanner.tex}, lines~400--403). The truncation estimates supply
a constant $C\ge1$ in $t\le Cn^{-500}$ chosen before the scale,
domain, and scan. By the spectral perturbation bound of~\cite[Proposition~4.5]{OpenAI2026AreaLaw},
a truncated ground-state energy $0\le\varepsilon\le n^{-1000}$ and pre-truncation
normalized spectral gap $g>0$ (where $g=\Delta/c_s$ when the physical
Hamiltonian has gap $\Delta$ and positive-constraint normalization $c_s$) give
\begin{align}
    t\le\sqrt{2\varepsilon/g},\qquad
    \varepsilon\le n^{-1000},\qquad g>0,
    \label{eq:arealaw2d_scanner_inputs_perturbation_bound}
\end{align}
and (\ref{eq:arealaw2d_scanner_inputs_perturbation_bound}) implies
$t\le\sqrt{2/g}\,n^{-500}$, so choosing $C\ge\max\{1,\sqrt{2/g}\}$
gives the required trace-distance bound.
For any fixed $A,C\ge0$ and local dimension $q\ge1$, every $t\le An^{-500}$
satisfies, for all sufficiently large $n$,
\begin{align}
    C\sqrt t(1+Cn^2\log q)\le n^{-247}.
    \label{eq:arealaw2d_scanner_inputs_continuity_cost}
\end{align}
Setting $A=C$ in (\ref{eq:arealaw2d_scanner_inputs_continuity_cost}) yields the
shell-entropy continuity bound used in the scanner deduction.

\section{Steps 2 and 3: Scan histories, sampling, and replica rounds}

\paragraph{Classical scan histories, sampling, and bad-history tails (Lemma~9.1).}
The finite history process of independent uniform offsets, deterministic row
fills, and labelled radius-$r_0$ charge choices is constructed on the induced
lattice graph. Under the row bound, bounded anchor multiplicity $\mu$, and cut
clearance, every splitting radius-$r_0$ ball on a good history has a unique
charge slot, chosen with conditional probability
$1/(2\lceil C_1nD\rceil)\ge[2(C_1+1)nD]^{-1}$, and moves the unassigned part of
its support.\footnote{Splitting-ball sampling:
    \leanid{TNLean.PEPS.AreaLaw.Scan.CollarScan.good_split_sampling_domainGraph}.}
Because the old middle lies inside the truncation set, every splitting
designated support coincides with its radius-$r_0$ ball, so the same expected
lower bound holds for designated-support incidences.\footnote{Designated-support
    sampling:
    \leanid{TNLean.PEPS.AreaLaw.Scan.CollarScan.good_designated_split_sampling_domainGraph},
    using \leanid{TNLean.PEPS.AreaLaw.truncationRadius_eq_of_mem_crossingLabels}.}
Moreover, the assigned-site lead and consumed-row coverage confine each
receiving--middle incidence (which contains the terminal metric split event of
\path{06-transport.tex}, lines~336--352) to an anchor interval of width
$D+3r_0+1$, bounding the total incidence weight across bands by $10D/m$ for
both fill and charge mixtures.\footnote{Mixture dilution bounds:
    \leanid{TNLean.PEPS.AreaLaw.Scan.CollarScan.designated_split_mixture_dilution_domainGraph}
    and \leanid{TNLean.PEPS.AreaLaw.Scan.CollarScan.fill_split_mixture_dilution_domainGraph}.}

For the bad-history probability, every bad endpoint traces back along a strictly
time-ordered path of $s$ selected charges with $D<4r_0s$ inside the last
$4n(r_0s+1)+1$ steps. A union bound over prefix times, sides, bands, and compact
endpoints gives the all-time bad-history probability bound
\begin{align}
    2K(|T|+nL)(N+1)
    \sum_{\substack{0\le s\le N \\D<4r_0s}}\rho^s,
    \qquad
    \rho=\frac{9e\,nr_0\mu(1+2(2r_0)(2r_0+1))}{2\lceil C_1nD\rceil}.
    \label{eq:arealaw2d_scanner_inputs_bad_tail}
\end{align}
Under the paper's rounded scale exponents, the geometric series in
(\ref{eq:arealaw2d_scanner_inputs_bad_tail}) decays faster than any inverse
power of $n$, and in particular is at most $n^{-200}$ for all sufficiently
large $n$.\footnote{Bad-history tail bounds:
    \leanid{TNLean.PEPS.AreaLaw.Scan.CollarScan.sum_historyWeight_bad_through_geometric_domainGraph_le}
    and \leanid{TNLean.PEPS.AreaLaw.Scan.eventually_bad_history_probability_le}.}

\paragraph{Band nesting, cut bounds, and replica transport trees (Proposition~7.4).}
In both completed ($k\le nm$) and pre-charge ($k+1\le nm$) statuses under
$D\ge1$, $r_0\le D$, $4D\le m$, and $8Km\le L$, the regions $U$ and $UY$ differ
from deterministic whole-row prefixes by at most $3nD$ sites, and the number of
crossing lattice edges is at most $20nD$ for $U$ and $UY$, $40nD$ for $Y$, and
$16nD$ for $P_0$ and $V$.\footnote{Prefix approximations and cut-edge bounds:
    \leanid{TNLean.PEPS.AreaLaw.Scan.CollarScan.state_prefix_approximation_domainGraph},
    \leanid{TNLean.PEPS.AreaLaw.Scan.CollarScan.oldChargeState_prefix_approximation_domainGraph},
    \leanid{TNLean.PEPS.AreaLaw.Scan.CollarScan.state_cut_bounds_domainGraph}, and
    \leanid{TNLean.PEPS.AreaLaw.Scan.CollarScan.oldChargeState_cut_bounds_domainGraph}.}
Disjoint bands have nested physical regions, which implies that their band
metrics commute on the symmetric replica subspace across all independent
histories and completed/pre-charge leaf pairs.\footnote{Cross-band commutation:
    \leanid{TensorPower.ReplicaTransport.bandMetric_commute_on_symmetric_of_nested},
    \leanid{TNLean.PEPS.AreaLaw.Scan.CollarScan.quantumChargeMove_apply}, and
    \leanid{TNLean.PEPS.AreaLaw.Scan.CollarScan.chargeTransportData_crossBandCommute}.}

The recursive history mean trees have leaf weights $m^{-K}(2M)^{-Kk}$ matching
the uniform history law, and their interpolation endpoints match the adjacent
fill and charge roots.\footnote{History mean trees:
    \leanid{TNLean.PEPS.AreaLaw.Scan.historyMeanTree},
    \leanid{TNLean.PEPS.AreaLaw.Scan.historyMeanTree_weight},
    \leanid{TNLean.PEPS.AreaLaw.Scan.historyMeanTree_eval_succ},
    \leanid{TNLean.PEPS.AreaLaw.Scan.charge_interpRoot_zero}, and
    \leanid{TNLean.PEPS.AreaLaw.Scan.charge_interpRoot_one}.}
Integrating the pointwise conditional-entropy sampling inequality against the
old transported state shows that the expected entropy gain dominates the
good-history split defect with coefficient
$[2(C_1+1)nD]^{-1}$.\footnote{Integrated sampling bound:
    \leanid{TNLean.PEPS.AreaLaw.Scan.CollarScan.actualChargeEntropyDefect_le_entropyGain_domainGraph}.}
With the log-dimension cap $\ell=1+(2r_0+1)^2\log q$ on moved subsystems,
Proposition~7.4 applies to every fill round, every charge round, and the
combined family of $2nm$ rounds for any nonzero symmetric replica vector
$\mathrm{pre}$ and $a\ell\le c_0$, with an entropy remainder
$0\le\beta_k\le C_\beta\log(k+2)$ and an energy remainder tending to zero as
$k\to\infty$ when $\overline H_k\mathrm{pre}=E_0\mathrm{pre}$, and telescopes
over the $2nm$ rounds.\footnote{Single-round, all-round, and telescoped
    transport bounds:
    \leanid{TNLean.PEPS.AreaLaw.Scan.CollarScan.actualFill_transport_domainGraph},
    \leanid{TNLean.PEPS.AreaLaw.Scan.CollarScan.actualCharge_transport_domainGraph},
    \leanid{TNLean.PEPS.AreaLaw.Scan.CollarScan.actualRounds_transport_domainGraph},
    \leanid{TNLean.PEPS.AreaLaw.Scan.CollarScan.sum_actualRoundData_logNormSq_sub}, and
    \leanid{TNLean.PEPS.AreaLaw.Scan.CollarScan.actualRounds_entropy_integral_le_domainGraph}.}

\paragraph{Transported leaf measures and integrated charge defect.}
The leaf measures on the one-copy unit sphere (from QICLean revision
\path{acfb5c685e732f987a5a96a5c079ba7f821c4715}) have interior mass $1/2$ for
$a\ge0$ and any nonzero symmetric replica vector, and their integrals reproduce
the round's entropy gain and good-history charge defect.\footnote{Round measure
    specializations:
    \leanid{TNLean.PEPS.AreaLaw.Scan.transportScanRound},
    \leanid{TNLean.PEPS.AreaLaw.Scan.CollarScan.actualFillScanRound},
    \leanid{TNLean.PEPS.AreaLaw.Scan.CollarScan.actualChargeScanRound}, and
    \leanid{TNLean.PEPS.AreaLaw.Scan.CollarScan.actualChargeScanRound_chargeDefect}.}
With the coherent symbol $F_h(\theta)$ formed from the split-label entropy terms
$\eta_{i,(h,\mathrm{old})}(\theta)$, the charge defect at replica count $k$ is
\begin{align}
    F_h(\theta)     & =\sum_i\eta_{i,(h,\mathrm{old})}(\theta), &
    \mathcal Q_k(p) & =\sum_{h\text{ good}}w_hJ_{h,k,p}(F_h).
    \label{eq:arealaw2d_scanner_inputs_coherent_symbol}
\end{align}
Hence $\mathcal Q_k(p)$ is almost-everywhere measurable on $(0,1)$ and interval
integrable on $[0,1]$.\footnote{Charge-defect regularity:
    \leanid{TNLean.PEPS.AreaLaw.Scan.aestronglyMeasurable_transportScanRound_chargeDefect},
    \leanid{TNLean.PEPS.AreaLaw.Scan.intervalIntegrable_transportScanRound_chargeDefect},
    \leanid{TNLean.PEPS.AreaLaw.Scan.CollarScan.aestronglyMeasurable_actualFillScanRound_chargeDefect},
    \leanid{TNLean.PEPS.AreaLaw.Scan.CollarScan.intervalIntegrable_actualFillScanRound_chargeDefect},
    \leanid{TNLean.PEPS.AreaLaw.Scan.CollarScan.aestronglyMeasurable_actualChargeScanRound_chargeDefect}, and
    \leanid{TNLean.PEPS.AreaLaw.Scan.CollarScan.intervalIntegrable_actualChargeScanRound_chargeDefect}.}

Setting $\alpha=[2(C_1+1)nD]^{-1}$ and $\ell=1+(2r_0+1)^2\log q$, the
pointwise bound $\alpha\mathcal Q(p)\le\mathcal G(p)$ combines with
Proposition~7.4 on a single charge round: for any nonzero symmetric vector with
transported norm $N(p)$, $a>0$ with $a\ell\le c_0$, replica count $r\ge0$, and
$0\le p_0\le p_1\le1$,
\begin{align}
    ra\alpha\int_{p_0}^{p_1}\mathcal Q(p)\,dp
     & \le \log N^2(p_0)-\log N^2(p_1)
    +(p_1-p_0)\bigl(C_{\mathrm{ent}}raKa^{1/4}\ell^{e_{\mathrm{ent}}}+\beta_r\bigr).
    \label{eq:arealaw2d_scanner_inputs_charge_integral}
\end{align}
Summing over all $nm$ charge stages $s\in\{0,\dots,nm-1\}$ and using the
nonnegativity of the intervening fill-round entropy gains (by strong
subadditivity) yields the all-round integrated defect bound of
Proposition~9.2 of~\cite{OpenAI2026AreaLaw}: writing
$N^{\mathrm{fill}}_{s,k,\mathrm{pre}}$ for the filtered norm after $s$ charges,
\begin{align}
    ka\alpha\sum_{s=0}^{nm-1}\int_{p_0}^{p_1}
    \mathcal Q_{s,k,\mathrm{pre}}(p)\,dp
     & \le \log (N^{\mathrm{fill}}_{0,k,\mathrm{pre}})^2
    -\log (N^{\mathrm{fill}}_{nm,k,\mathrm{pre}})^2                      \\
     & \quad +2nm\bigl(C_{\mathrm{ent}}kaKa^{1/4}\ell^{e_{\mathrm{ent}}}
    +\beta_k\bigr).
    \label{eq:arealaw2d_scanner_inputs_charge_sum}
\end{align}
Both (\ref{eq:arealaw2d_scanner_inputs_charge_sum}) and the single-round bound
require only interval integrability of the scalar defect.\footnote{One-round
    and all-round charge integral bounds:
    \leanid{TNLean.PEPS.AreaLaw.Scan.CollarScan.actualCharge_integral_bound_domainGraph}
    and \leanid{TNLean.PEPS.AreaLaw.Scan.CollarScan.actualCharges_integral_bound_domainGraph}.}

\section{Step 3: Leaf masses and parameter regularity}

In Proposition~7.4 (\path{06-transport.tex}, lines~393--400), the transported
leaf states are normalized density matrices for interior parameters $0<p<1$, so
their combined Fourier--coherent measures have total mass $1/2$ on $(0,1)$,
whereas at $p=1$ (for old leaves, where $(1-p)w_h=0$) and $p=0$ (for new leaves,
where $pw_hq_{c\mid h}=0$) the leaf weights and transported measures vanish.
The mass-$1/2$ identity is used only for $0<p<1$ and is derived from these
normalized transported states. Since $\{0,1\}$ has Lebesgue measure zero,
interior bounds and almost-everywhere measurability suffice for integrating
over $[0,1]$.

Parameter integrability of the old-state Fourier--coherent integral is proved
from the transported expectation formula. For a continuous real coherent symbol
$F$ on the unit sphere with coherent average $A_F$, Hermiticity of
$Q_F=(A_F+A_F^\dagger)/2$ and positivity of the state $\rho$ give
\begin{align}
    \operatorname{Re}\operatorname{Tr}(\rho Q_F)
    =\operatorname{Re}\operatorname{Tr}(\rho A_F)
    =\mathcal H_\rho(F).
    \label{eq:arealaw2d_scanner_inputs_coherent_hermitian}
\end{align}
Measurability in $p$ and the pointwise bound $|J_h(F)|\le\|Q_F\|/2$ imply
interval integrability on $[0,1]$ for each leaf and hence for finite weighted
sums over good histories.\footnote{Integrability declarations:
    \leanid{TensorPower.ReplicaTransport.TransportData.intervalIntegrable_oldFourierCoherentIntegral},
    \leanid{TensorPower.ReplicaTransport.TransportData.intervalIntegrable_sum_oldFourierCoherentIntegral},
    and \leanid{TNLean.PEPS.AreaLaw.Scan.CollarScan.intervalIntegrable_actualChargeEntropyDefect}.}

\section{Step 4: Truncated ground state, marginal tails, and typical window}

\paragraph{Truncated ground state and status marginal moments.}
The log-moment generating function and spectral tail bounds
of~\cite[Lemma~3.1]{OpenAI2026AreaLaw} are formalized in
\path{TNLean/PEPS/AreaLaw/Scan/ActualMarginalTails.lean} for every regional
marginal of a unit ground vector $\widetilde\Omega$ of the truncated
Hamiltonian $\widetilde H=\sum_i\widetilde h_i$ satisfying
\begin{align}
    \|\widetilde\Omega\|         & =1,                 \\
    \widetilde H\widetilde\Omega & =e\widetilde\Omega, \\
    \widetilde H-eI              & \ge(g/2)
    (I-|\widetilde\Omega\rangle\langle\widetilde\Omega|),
    \label{eq:arealaw2d_scanner_inputs_marginal_gap}
\end{align}
where $g=\Delta/\max\{1,J(\|f_{1,\Delta/2}\|_1+1)\}$ is the normalized spectral
gap determined by the original gap $\Delta$, the term norm bound $J$, and the
$L^1$ norm of the quasi-local filter kernel $f_{1,\Delta/2}$
of~\cite[Lemma~4.2]{OpenAI2026AreaLaw} under the positive-constraint
normalization of~\cite[Section~4]{OpenAI2026AreaLaw}.

By Proposition~4.5, for $r_0=\lceil C_{\mathrm{tr}}(\log n)^2\rceil$ and a
compact scan neighborhood $S_0$ with $|S_0|\le C_0n^2$, there exists a pair
$(e,\widetilde\Omega)$ satisfying (\ref{eq:arealaw2d_scanner_inputs_marginal_gap})
that gives the truncated Hamiltonian $\widetilde H$ in
(\ref{eq:arealaw2d_scanner_inputs_marginal_gap}) and obeys, with
$\epsilon_n=\min\{n^{-1000},g/4\}$ and some phase $\theta\in\mathbb R$,
\begin{align}
    0 & \le e\le\epsilon_n,      \\
    \|\widetilde\Omega-\mathrm e^{i\theta}\Omega\|
      & \le2\sqrt{\epsilon_n/g}, \\
    \frac12\bigl\||\widetilde\Omega\rangle\langle\widetilde\Omega|
    -|\Omega\rangle\langle\Omega|\bigr\|_1
      & \le\sqrt{2\epsilon_n/g}.
    \label{eq:arealaw2d_scanner_inputs_marginal_closeness}
\end{align}
Across all five status regions $XU,V,Y,U,UY$ of every completed and pre-charge
step, the crossing capacity satisfies $1\le\mathcal B\le C_BnD(\log n)^{12}$,
and the marginal log-moment bound holds on $|u|\le c_0/\sqrt{\mathcal B}$ with
$C_0=512\mathrm e(2/g)$ and $c_0=[32\sqrt{1+2/g}]^{-1}$ while retaining
(\ref{eq:arealaw2d_scanner_inputs_marginal_closeness}).\footnote{Truncated
    marginal moment and tail declarations:
    \leanid{TNLean.PEPS.AreaLaw.Scan.CollarScan.exists_truncated_reducedState_moment_tail_bounds},
    \leanid{TNLean.PEPS.AreaLaw.Scan.CollarScan.exists_truncated_statusMarginal_moment_bounds},
    and \leanid{TNLean.PEPS.AreaLaw.Scan.CollarScan.exists_source_statusMarginal_moment_bounds}.}

\paragraph{Target typical window and final collar parameters.}
In \path{TNLean/PEPS/AreaLaw/Scan/TargetTypicalWindow.lean}, under the paper's
target geometry~\cite{OpenAI2026AreaLaw} (\path{08-scanner.tex}, lines~21--48,
243--257, and 400--414)---a nonempty ambient target $T$ with $|T|\le C_Tn^2$,
at most $n$ sites per ambient layer up to $L=\lfloor n^{1-\ell_T}\rfloor$, and
clearance $\operatorname{dist}_\infty(T,Z)>2L+10r_0$ from the endpoint set $Z$
of $\partial_\Lambda A$---the target $X=A\cap T$ and collar neighborhood
$S_0=\{x\in A:\operatorname{dist}_\infty(x,T)\le L\}$ satisfy\footnote{Target
    geometry declarations:
    \leanid{TNLean.PEPS.AreaLaw.Scan.ScannerExponents.one_le_L_le} and
    \leanid{TNLean.PEPS.AreaLaw.Scan.CollarScan.source_target_geometry}.}
\begin{align}
    |S_0|\le(C_T+1)n^2,\qquad X\subseteq S_0,\qquad
    |\partial_\Lambda X|\le4n.
    \label{eq:arealaw2d_scanner_inputs_target_geometry}
\end{align}
Because $|\partial_\Lambda X|\le 4n$, the crossing capacity $\mathcal B_X$ of
$X$ has no factor of $D$ and satisfies\footnote{Target crossing capacity:
    \leanid{TNLean.PEPS.AreaLaw.one_le_cutBudget} and
    \leanid{TNLean.PEPS.AreaLaw.Scan.CollarScan.exists_truncated_target_typical_window}.}
\begin{align}
    1\le\mathcal B_X\le C_Bn(\log n)^{12}.
    \label{eq:arealaw2d_scanner_inputs_target_budget}
\end{align}
For the same truncated ground state $(e,\widetilde\Omega)$ satisfying
(\ref{eq:arealaw2d_scanner_inputs_marginal_gap}) and
(\ref{eq:arealaw2d_scanner_inputs_marginal_closeness}), the spectral tail of the
reduced state on $X$ outside the surprisal window of width $w=n^{3/5}$ around
$s_X=S_{\widetilde\Omega}(X)$~\cite{OpenAI2026AreaLaw}
(\path{07-comparators.tex}, lines~37--53) is at most $n^{-100}$ for $n\ge N$,
so the typical eigenvalue set is nonempty and its total mass $z$ and each
typical eigenvalue $p_i$ satisfy\footnote{Target typical window:
    \leanid{TNLean.PEPS.AreaLaw.Scan.CollarScan.exists_truncated_target_typical_window}.}
\begin{align}
    1-n^{-100}\le z\le1,\qquad
    \mathrm e^{-s_X-w}\le p_i\le\mathrm e^{-s_X+w}.
    \label{eq:arealaw2d_scanner_inputs_target_typical_mass}
\end{align}
Finally, for weights $0\le W\le n^\omega$ and crossing capacities
$1\le B\le C_BnD(\log n)^{12}$, the exponent inequality
\begin{align}
    \omega+(1+\kappa)/2 & <1-\ell_{\rm s}-\mu
    \label{eq:arealaw2d_moment_margin}
\end{align}
ensures that $Aa\le c/\sqrt B$ (where $a=W/K$) holds for all sufficiently large
$n$; while condition (\ref{eq:arealaw2d_moment_margin}) is not used in the
bootstrap step, the final collar parameters of~\cite[Section~9, lines~839--865]{OpenAI2026AreaLaw}
satisfy it with margin $499879/2000000$ at $W=n^{1/20000}$.\footnote{Moment-radius
    scale bounds:
    \leanid{TNLean.PEPS.AreaLaw.Scan.ScannerExponents.exists_a_le_moment_radius}
    and \leanid{TNLean.PEPS.AreaLaw.Scan.exists_final_a_le_moment_radius}.}

\section{Step 5: Charge-round energy and H\"older accounting}

For the finite-leaf energy error $\mathcal E_k(p)$ of~\cite{OpenAI2026AreaLaw}
(\path{08-scanner.tex}, lines~416--440 and~492--511) in charge round $s+1\le nm$
under $n\ge2$, $1\le D$, $r_0\le D$, $4(D+2r_0)\le m$, $8Km\le L$, and anchor
multiplicity at most $\mathfrak m$, set $A=10KnD\mathfrak m$,
$\ell=1+(2r_0+1)^2\log q$, and let $b=\sum_{h\text{ not good}}w_h$ be the
bad-old-history probability. For any nonzero symmetric replica vector $\psi_k$
and $0<p<1$, H\"older's inequality and the $10D/m$ dilution bound yield
\begin{align}
    \mathcal E_k(p) & \le\frac{10D}{m}
    \left((1-p)A^{7/8}\mathcal Q_k(p)^{1/8}
    +\frac{A(2\ell)^{1/8}}2\bigl((1-p)b+p\bigr)\right).
    \label{eq:arealaw2d_scanner_inputs_charge_energy}
\end{align}
At the scanner exponents $\ell_{\mathrm{sc}}=10^{-5}$, $\kappa=10^{-6}$, and
$\mu_{\mathrm{sc}}=1/4$, the bad-history tail bound replaces $b$ in
(\ref{eq:arealaw2d_scanner_inputs_charge_energy}) by $n^{-P}$ for any fixed
$P$ and all sufficiently large $n$.\footnote{Charge-round energy bounds:
    \leanid{TNLean.PEPS.AreaLaw.Scan.CollarScan.actualCharge_energyError_le_domainGraph},
    \leanid{TNLean.PEPS.AreaLaw.Scan.exists_badOldHistory_threshold}, and
    \leanid{TNLean.PEPS.AreaLaw.Scan.exists_actualCharge_energy_threshold}.}
Absorbing the transport prefactors $|C_{\mathrm{en}}|a^2\ell^{e_{\mathrm{en}}}$
with $a=W/K$ into a power of $\log n$ then gives, for all $n\ge N$,
\begin{align}
    |C_{\mathrm{en}}|a^2\ell^{e_{\mathrm{en}}}\mathcal E_k(p)
     & \le C W^2 n^{\ell_{\mathrm{sc}}}D^2(\log n)^{C_l}
    \left((1-p)\left(\frac{\mathcal Q_k(p)}{KnD}\right)^{1/8}
    +(1-p)n^{-P}+p\right),
    \label{eq:arealaw2d_scanner_inputs_charge_coefficients}
\end{align}
matching the energy-error functional of lines~507--511 prior to round
selection.\footnote{Defect nonnegativity and normalized energy coefficients:
    \leanid{TNLean.PEPS.AreaLaw.Scan.CollarScan.actualChargeEntropyDefect_nonneg} and
    \leanid{TNLean.PEPS.AreaLaw.Scan.exists_actualCharge_energy_coefficients}.}

\section{Effect on the source statement}

All four comparisons refer to the same selected charge and interpolation
parameter, with the same truncated vector, typical set and label sequence.
The sharp comparisons retain precisely their marginal smallness
antecedent. No transport or comparison certificate is an additional
hypothesis of the physical scanner theorem.

The retained local adjustment distinguishes the input entropy coefficient
from output coefficients enlarged after it is fixed. It does not change
the powers of $n$, the rounded scales, the admissible exponents, the total
weight condition or the order of the physical and replica quantifiers.
The identical-coefficient reading of the printed symbol $C_e$ is not
claimed. The explicit marginal logarithmic exponent $12$ and the
continuity bound $n^{-247}$ are quantitative choices used to absorb the
errors into the allowed output constants.

\bibliographystyle{alpha}
\bibliography{references}

\end{document}
